\documentclass{article}
\usepackage{graphicx} % Required for inserting images

\title{finiteness of G-simulation}
\author{chen zhao}
\date{January 2026}

\begin{document}

\maketitle

% ============================================================
% Meeting-style explanation for advisor discussion
% PDTA termination with G-simulation (diagonal constraints)
% Based strictly on:
%   [PDTA] https://arxiv.org/pdf/2105.13683
%   [GAbs] https://arxiv.org/pdf/1904.08590
% ============================================================

\section*{Termination of PDTA Reachability with G-simulation}

\paragraph{Goal of the discussion.}
I want to explain why the PDTA reachability algorithm
terminates when we replace LU-extrapolation by
\emph{G-simulation} in order to support diagonal constraints.

The key point is to show that the simulation relation used for
subsumption induces a \emph{finite symbolic state space},
so that the saturation algorithm must stop.

\vspace{1ex}
The argument follows exactly the structure of:
\begin{itemize}
  \item the PDTA termination theorem (Theorem~3 in~\cite{PDTA}), and
  \item the finiteness proof of G-simulation in~\cite{GAbs}.
\end{itemize}

% ------------------------------------------------------------

\subsection*{1. What ``finite'' means in the PDTA algorithm}

In the PDTA paper, symbolic states are pairs $(q,Z)$
(control state + zone).
The algorithm explores these states using inference rules.

There are two critical finiteness requirements:

\begin{enumerate}
  \item The set of \emph{roots} $S^\star$ must be finite.
  \item For each root $(q,Z)$, the local exploration set $S(q,Z)$
        must be finite.
\end{enumerate}

This is exactly Theorem~3(1) in~\cite{PDTA}.

The PDTA proof shows:
\begin{itemize}
  \item Root finiteness follows if the simulation equivalence
        $\sim$ has \emph{finite index} (called \emph{strong finiteness}).
  \item Local finiteness follows if the simulation preorder
        $\preceq$ is a \emph{finite simulation}.
\end{itemize}

So the whole question reduces to:
\begin{quote}
\emph{Why is G-simulation finite (and strongly finite) when diagonal
constraints are allowed?}
\end{quote}

% ------------------------------------------------------------

\subsection*{2. High-level idea of G-simulation}

The difficulty is that diagonal constraints break classical
LU-extrapolation: $4_{LU}$ is no longer a simulation.

The solution proposed in~\cite{GAbs} is:
\begin{itemize}
  \item Use \emph{state-based guard sets} $G(q)$,
  \item Define a simulation that combines:
    \begin{enumerate}
      \item LU-style simulation for non-diagonal constraints, and
      \item implication checking for diagonal constraints.
    \end{enumerate}
\end{itemize}

This gives a simulation that is:
\begin{itemize}
  \item sound for timed automata with diagonals, and
  \item still finite, hence usable for pruning.
\end{itemize}

% ------------------------------------------------------------

\subsection*{3. What ``finite simulation'' means (key definition)}

Following~\cite{GAbs}, a simulation $\triangleleft$
is called \emph{finite} if the mapping
\[
Z \longmapsto \downarrow Z
\quad\text{with}\quad
\downarrow Z = \{v \mid \exists v' \in Z : v \triangleleft v'\}
\]
has \emph{finite range}.

Intuitively:
\begin{quote}
There are only finitely many distinct ``behavioural shapes''
of zones under subsumption.
\end{quote}

This is exactly what ensures that forward reachability
with subsumption cannot generate infinitely many new nodes.

% ------------------------------------------------------------

\subsection*{4. Detailed Step 1: Why G-simulation is finite}

This is the technical core.

\subsubsection*{4.1 Decomposition of the simulation}

The valuation-level G-simulation can be written as:
\[
v \sqsubseteq^{LU}_G v'
\iff
\underbrace{v\ 4_{LU(G)}\ v'}_{\text{non-diagonal part}}
\ \wedge\
\underbrace{
\forall \varphi \in G^+:\ (v\models\varphi \Rightarrow v'\models\varphi)
}_{\text{diagonal part}}
\]

where:
\begin{itemize}
  \item $G^+$ is the finite set of diagonal constraints in $G(q)$,
  \item $LU(G)$ is extracted from $G(q)$.
\end{itemize}

So finiteness is proved by analyzing these two parts separately.

\subsubsection*{4.2 Non-diagonal part: inherited LU finiteness}

The relation $4_{LU(G)}$ is exactly the classical LU-simulation.

By Theorem~4 recalled in~\cite{GAbs}:
\begin{quote}
$4_{LU(G)}$ is a finite simulation.
\end{quote}

So the non-diagonal part already has only finitely many downsets.

\subsubsection*{4.3 Diagonal part: finite Boolean abstraction}

This is the crucial diagonal argument.

Because $G^+$ is finite, each valuation $v$ induces a Boolean vector:
\[
\chi(v) \in \{0,1\}^{|G^+|}
\quad\text{where}\quad
\chi(v)_\varphi = 1 \iff v\models\varphi.
\]

Hence:
\begin{itemize}
  \item The valuation space is partitioned into at most
        $2^{|G^+|}$ diagonal classes.
  \item The diagonal implication check depends only on these classes.
\end{itemize}

Therefore, the diagonal part of the simulation also has
only finitely many possible downsets.

\subsubsection*{4.4 Combining both parts}

The full G-simulation is the conjunction of:
\begin{itemize}
  \item a finite LU-simulation, and
  \item a finite diagonal-implication relation.
\end{itemize}

Thus the overall downset mapping has finite range.

\begin{quote}
\emph{Conclusion: G-simulation is a finite simulation}
(Lemma~15 in~\cite{GAbs}).
\end{quote}

This implies that zone graph exploration with G-simulation
is terminating at the timed-automaton level.

\subsubsection*{4.5 Current implementation correspondence}

In the current prototype implementation, the theory objects are reflected
by a cached state summary attached to the symbolic control part:
\begin{itemize}
  \item the collected diagonal constraints correspond to $G^+(q)$,
  \item the simple constraints are projected to local arrays $L,U$,
        which implement the current approximation of $LU(G(q))$,
  \item subsumption first checks $aLU$-inclusion using these local
        bounds, and then checks diagonal implication using $G^+(q)$.
\end{itemize}

So the implementation now follows the intended decomposition
\[
\text{LU}(G(q)) \quad + \quad \text{diagonal implication over } G^+(q),
\]
although a fully explicit Boolean abstraction layer is still future work.

% ------------------------------------------------------------

\subsection*{5. Lifting finiteness to PDTA}

Now we plug this result into the PDTA framework.

\subsubsection*{5.1 Root finiteness}

In PDTA, only \textsc{Push} creates new roots.
A push is added only if there is no equivalent root under~$\sim$.

The PDTA paper explicitly states that G-simulation belongs to the class of
\emph{strongly finite simulations}.

\medskip
\noindent
\emph{Hence, only finitely many roots can be created.}


\subsubsection*{5.2 Local finiteness}

Inside each context $S(q,Z)$:
\begin{itemize}
  \item Internal and Pop rules add nodes only if they are
        not subsumed.
  \item Since G-simulation is finite, only finitely many
        non-subsumed nodes can appear.
\end{itemize}

So each $S(q,Z)$ is finite.

% ------------------------------------------------------------

\subsection*{6. Final conclusion}

Combining everything:

\begin{itemize}
  \item G-simulation yields a finite subsumption relation
        even with diagonal constraints.
  \item Root creation is finite (strong finiteness).
  \item Local exploration is finite (finite simulation).
\end{itemize}

\begin{quote}
\textbf{Therefore, the PDTA reachability algorithm
terminates when G-simulation is used to support diagonal constraints.}
\end{quote}

This is exactly the finiteness argument required by
Theorem~3 of~\cite{PDTA}, instantiated with the
G-simulation construction of~\cite{GAbs}.

% ------------------------------------------------------------
% Bibliography (explicit papers instead of placeholders)
% ------------------------------------------------------------

\begin{thebibliography}{9}

\bibitem{PDTA}
Akshay, P., Gastin, P., Prakash, N.
\newblock \emph{Fast Zone-Based Algorithms for Reachability in Pushdown Timed Automata}.
\newblock arXiv:2105.13683.
\newblock \url{https://arxiv.org/pdf/2105.13683}

\bibitem{GAbs}
Gastin, P., Mukherjee, K. N., Srivathsan, B.
\newblock \emph{Fast Algorithms for Handling Diagonal Constraints in Timed Automata}.
\newblock arXiv:1904.08590.
\newblock \url{https://arxiv.org/pdf/1904.08590}

\end{thebibliography}

\end{document}
